\documentclass[12pt]{article}
\usepackage[T1]{fontenc}
\usepackage[utf8]{inputenc}
\usepackage{lmodern}
\usepackage{textcomp}
\usepackage{enumitem}
\usepackage{geometry}
\geometry{margin=1in}
\begin{document}
\begin{center}
Answer Key - Chemistry - Undergraduate Level Assessment 14 \
\small (Answers are ordered exactly as in the question paper.)
\end{center}

\section*{Multiple Choice}
\begin{enumerate}[label=\arabic*.]
\item \textbf{Answer:} A
\\ \textit{Options:} A) 17 O; B) 19 F; C) 29 Si; D) 31 P
\\ \textit{Note:} The correct answer is 17 O because its nuclear magnetic resonance (NMR) frequency of 115.5 MHz at a 20.0 T magnetic field is consistent with its gyromagnetic ratio, which determines the resonance frequency of a nuclide in a given magnetic field strength.
\item \textbf{Answer:} A
\\ \textit{Options:} A) Neutrons; B) Alpha particles; C) Beta particles; D) X-rays
\\ \textit{Note:} Cobalt-60 is produced from the neutron bombardment of cobalt-59 through a nuclear reaction, where a neutron is absorbed by cobalt-59, resulting in the formation of cobalt-60, which is a radioactive isotope used in cancer radiation therapy.
\item \textbf{Answer:} D
\\ \textit{Options:} A) I only; B) II only; C) III only; D) I and III only
\\ \textit{Note:} The correct answer is I and III only because infrared-active modes require a change in the dipole moment during vibration, which occurs in the bending and asymmetric stretching modes of carbon dioxide, while the symmetric stretching mode does not produce a change in dipole moment and is therefore infrared-inactive.
\item \textbf{Answer:} C
\\ \textit{Options:} A) 0.5 T; B) 1.2 T; C) 2.9 T; D) 100 T
\\ \textit{Note:} The correct answer is 2.9 T because it is derived from the relationship between the magnetic field strength, polarization, and the temperature for a paramagnetic substance, using the appropriate equations that relate these variables.
\item \textbf{Answer:} B
\\ \textit{Options:} A) 0.0471 Hz; B) 21.2 Hz; C) 42.4 Hz; D) 66.7 Hz
\\ \textit{Note:} The linewidth in NMR is inversely related to the T$_{2}$ relaxation time, and can be calculated using the formula linewidth (Hz) = 1/(2πT 2), resulting in a linewidth of approximately 21.2 Hz for a T$_{2}$ of 15 ms.
\end{enumerate}

\section*{True / False}
\begin{enumerate}[label=\arabic*.]
\setcounter{enumi}{5}
\item \textbf{Answer:} False
\\ \textit{Options:} A) True; B) False
\\ \textit{Note:} The statement is false because, unlike alkali metals, lanthanides do not react vigorously with aqueous acids to liberate hydrogen gas; their reactions are generally less intense and depend on specific conditions.
\item \textbf{Answer:} True
\\ \textit{Options:} A) True; B) False
\\ \textit{Note:} The NMR frequency of 31 P in a 20.0 T magnetic field is true because the Larmor frequency, which is directly proportional to the strength of the magnetic field and the gyromagnetic ratio of 31 P, results in a calculated frequency of approximately 345.0 MHz.
\item \textbf{Answer:} False
\\ \textit{Options:} A) True; B) False
\\ \textit{Note:} HCl is an acid that donates protons, while NaOH is a base that accepts protons, and they do not form a conjugate acid-base pair because they do not convert into each other through proton transfer.
\item \textbf{Answer:} True
\\ \textit{Options:} A) True; B) False
\\ \textit{Note:} The statement is true because hexane has a symmetrical structure that leads to distinct environments for its carbon atoms, resulting in three unique chemical shifts in its 13 C NMR spectrum.
\item \textbf{Answer:} True
\\ \textit{Options:} A) True; B) False
\\ \textit{Note:} At standard temperature and pressure, bromine (Br$_{2}$) is the only non-metallic element that is a liquid, characterized by its red-brown color and volatility.
\end{enumerate}

\section*{Long Form Response}
\begin{enumerate}[label=\arabic*.]
\setcounter{enumi}{10}
\item \textbf{Answer:} Doubling the volume of a buffer solution composed of equal concentrations of a weak acid and its conjugate base has little effect on its pH. This is because the pH of a buffer solution is determined by the ratio of the concentrations of the weak acid and its conjugate base, as described by the Henderson-Hasselbalch equation, which states that pH = pKa + log([A-]/[HA]). When the volume is doubled, the concentrations of both the acid and the base are halved, but their ratio remains unchanged. Therefore, despite the increase in volume, the pH remains stable, illustrating the buffer solution's ability to resist changes in pH when diluted.
\\ \textit{Note:} correct because the pH of a buffer solution is determined by the ratio of the concentrations of its weak acid and conjugate base, which remains constant when the volume is doubled, leading to no significant change in pH despite the halving of individual concentrations.
\item \textbf{Answer:} Lattice enthalpy in ionic compounds is influenced by several key factors, including the charges of the ions and the distance between them. Higher charges lead to stronger electrostatic forces, while smaller ionic radii decrease the distance between ions, further increasing lattice enthalpy. Magnesium oxide (MgO) exhibits the greatest lattice enthalpy among common ionic substances primarily due to the high charges of its constituent ions, $Mg^2$⁺ and $O^2$⁻, which generate particularly strong electrostatic attractions. Additionally, the small size of these ions allows for closer packing in the crystal lattice, enhancing the overall stability of the compound. As a result, the combination of high ionic charges and small ionic radii in MgO contributes to its exceptionally high lattice enthalpy.
\\ \textit{Note:} correct because it accurately identifies that the lattice enthalpy is determined by the ionic charges and distances, and it specifically highlights that magnesium oxide has high charges on its ions, leading to stronger electrostatic interactions and hence the greatest lattice enthalpy among common ionic compounds.
\end{enumerate}

\vfill
\noindent\textit{End of Answer Key}
\end{document}